\chapter{From learning server to production ecosystem}
\chapterintro{Building the stack makes mature web infrastructure less magical. It should also make us humble: production systems solve a broader, continuously changing problem.}

\begin{objectives}
\begin{itemize}
  \item Explain the roles of WSGI, ASGI, production servers, frameworks, and reverse proxies.
  \item Identify the authentication, deployment, observability, and operations work omitted here.
  \item Transfer low-level lessons to everyday backend development.
  \item Complete a final assessment and choose safe next projects.
\end{itemize}
\end{objectives}

\section{What frameworks no longer hide}

When a mature framework gives a request object, you now know the work behind it:
stream reads, framing, syntax, field normalization, limits, target parsing, routing,
content decoding, exception boundaries, and response serialization. When a
production server exposes worker and timeout settings, you can connect them to
capacity and slow-client behavior rather than treating them as arbitrary knobs.

This knowledge changes debugging. A 400 before application logs suggests an
earlier parser or proxy. A 405 differs from a missing route. A hanging client can
indicate missing length or an incomplete stream. A database error belongs after
valid HTTP and domain input.

\section{WSGI and ASGI}

WSGI defines a synchronous interface between Python web servers and applications.
The server handles network protocol work and calls an application with request
environment data; the application returns status, fields, and body iterable.
It enabled portable synchronous applications across many servers.

ASGI generalizes the boundary for asynchronous and long-lived interactions. An
ASGI application receives a connection scope and exchanges event messages through
\code{receive} and \code{send} callables. It supports HTTP, WebSocket, lifespan
events, and asynchronous application work.

Neither interface is ``HTTP itself.'' They are Python application/server contracts
above protocol parsing. Our \code{Application.handle(Request) -> Response} is a
smaller educational boundary with similar separation but no compatibility claim.

\section{A typical production path}

A deployment might place these responsibilities in layers:

\begin{enumerate}
  \item edge network and firewall absorb or filter unwanted traffic;
  \item load balancer or reverse proxy terminates TLS, enforces limits, and routes;
  \item production WSGI/ASGI server manages processes, workers, timeouts, and protocol details;
  \item framework supplies routing, validation integration, middleware, and application tools;
  \item application implements resource and authorization rules;
  \item database/cache/message systems provide durable or shared state;
  \item logs, metrics, traces, alerts, deployment, backups, and incident processes operate it.
\end{enumerate}

Each boundary has timeout and size policy. If an outer proxy allows a 100 MiB
body while an inner service allows 1 MiB, rejection occurs later after wasting
bandwidth. If an inner timeout exceeds every outer timeout, work can continue
after the client has already received a gateway error. Production configuration
is an end-to-end design.

\section{Authentication and authorization}

Authentication establishes a principal; authorization decides what that principal
may do. Adding users safely requires choices about password hashing, multi-factor
authentication, account recovery, session or token lifetime, revocation, CSRF,
credential storage, TLS, rate limiting, audit, privacy, and authorization at
resource boundaries.

Use established framework and identity-system capabilities. Do not hash passwords
with a fast general-purpose hash, invent token signatures, store secrets in source,
or trust identity fields from untrusted proxies. Even a correct cryptographic
primitive can be assembled into an insecure authentication system.

\section{Observability and operations}

Access logs answer individual request questions. Metrics reveal rates, latency
distributions, saturation, queue depth, error classes, and dependency health.
Traces connect work across services. Alerts should correspond to user impact and
actionable operator decisions.

Operations also includes reproducible builds, dependency updates, vulnerability
response, configuration management, least privilege, secret rotation, database
migrations, backups, restore tests, capacity planning, rolling deploys, rollback,
and incident learning. A server is not production-ready merely because it answers
requests quickly in a benchmark.

\section{Safe next projects}

Good continuations preserve one learning question at a time:

\begin{itemize}
  \item implement chunked request decoding with a written grammar and adversarial tests;
  \item add TLS only through Python's established \code{ssl} module and local certificates,
        studying verification rather than implementing cryptography;
  \item adapt the Task application to WSGI, comparing boundaries with tinyhttp;
  \item adapt it to ASGI and measure slow concurrent clients;
  \item add conditional updates using ETag and \code{If-Match};
  \item replace offset pagination with a documented stable cursor;
  \item run the API behind a local reverse proxy and trace forwarded metadata trust.
\end{itemize}

Do not expose \code{tinyhttp} publicly as the next project. The goal is transfer
of understanding to maintained tools, not a new production server ecosystem.

\section{Final assessment}

The intern should complete four demonstrations without copying prose:

\begin{enumerate}
  \item On a whiteboard, trace an HTTPS Task POST from URL to DNS, TCP/TLS, HTTP,
        parser, router, validation, transaction, and response.
  \item Starting from a blank file, create a one-response socket server and then
        describe every missing property required for the final subset.
  \item Diagnose five supplied failures using wire bytes, logs, status, and tests,
        locating the owning layer before changing code.
  \item Extend the Task API with one read-only aggregate endpoint, including
        contract, repository method, handler, unit tests, network test, and docs.
\end{enumerate}

Assessment emphasizes explanation and evidence. A working patch without a clear
boundary model is incomplete; an elegant explanation without tested bytes is
also incomplete.

\section{The central lesson}

An HTTP backend is a sequence of contracts. Bytes become a framed message;
syntax becomes typed internal state; paths and methods become a resource
operation; untrusted values become validated data; a transaction becomes durable
state; a response value becomes bytes again. Reliability comes from making each
contract explicit, bounded, testable, and honest about what it does not support.

\summarytitle

WSGI, ASGI, mature servers, frameworks, proxies, identity systems, and operational
platforms divide responsibilities that our small project exposes directly.
Low-level understanding improves everyday design and debugging, while production
readiness requires maintained implementations and organizational processes far
beyond a successful local API.

\reviewtitle
\begin{enumerate}
  \item At what layer do WSGI and ASGI sit relative to HTTP parsing?
  \item Why must timeout policy be coordinated across proxies and applications?
  \item What work turns authentication from a token check into a system?
  \item Which tinyhttp lessons transfer directly to a mature framework?
\end{enumerate}

\begin{exercise}[Production comparison]
Choose one maintained Python framework and production server. Map their request,
response, routing, middleware, worker, limit, timeout, and shutdown concepts to
tinyhttp. Cite official documentation and identify capabilities with no tinyhttp
equivalent.
\end{exercise}

\begin{exercise}[Capstone retrospective]
Write a two-page engineering note: three assumptions you held before the book,
three experiments that changed your model, the most security-sensitive parser
decision, and the production boundary you would defend most strongly in review.
\end{exercise}

